\documentclass[9pt]{beamer}
\usepackage[utf8]{inputenc}
% \usepackage{enumitem}  % Conflicts with beamer's enumerate
\usepackage{booktabs}
\usepackage{fontspec}
% \usepackage{unicode-math}

\setmainfont{Times New Roman}
% \setmathfont{TeX Gyre Termes Math}  % Commented out - font not installed

\usetheme{Madrid}
\usecolortheme{default}

\usepackage[
  backend=biber,
  style=apa,
  citestyle=apa,
  natbib=false,
  doi=false,
  isbn=false,
  url=false
]{biblatex}

\addbibresource{delivery_emission.bib}

\usepackage{caption}
\setbeamertemplate{caption}[numbered]

%------------------------------------------------------------
%This block of code defines the information to appear in the
%Title page
\title[Drones for Food Delivery] %optional
{Drones vs. E-bikes: Comparing Operational Emissions in Urban Food Delivery in China}

% \subtitle{A short story}

\author[Jingxian Lin] % (optional)
{Jingxian Lin}

\institute[HKU] % (optional)
{
  Department of Civil Engineering\\
  The University of Hong Kong
}

\date[\today] % (optional)
{Seminar for Confirmation of PhD Candidature}

%End of title page configuration block
%------------------------------------------------------------



%------------------------------------------------------------
%The next block of commands puts the table of contents at the 
%beginning of each section and highlights the current section:

\AtBeginSection[]
{
  \begin{frame}
    \frametitle{Table of Contents}
    \tableofcontents[currentsection]
  \end{frame}
}
%------------------------------------------------------------


\begin{document}

%The next statement creates the title page.
\frame{\titlepage}


%---------------------------------------------------------
%This block of code is for the table of contents after
%the title page
\begin{frame}
\frametitle{Table of Contents}
\tableofcontents
\end{frame}
%---------------------------------------------------------


\section{Introduction}


\begin{frame}{Food Delivery in China}
\begin{figure}
    \centering
    \includegraphics[width=0.8\linewidth]{../figures/market/china_food_delivery_market_size.pdf}
    \caption{\textbf{China's online food delivery market size (2017–2027)}, showing historical values and forecasts \parencite{statista2024chinafooddelivery}. The market has grown from approximately 300 billion Yuan in 2017 to over 1 trillion Yuan by 2023, with projections suggesting continued expansion.}
    \label{fig:market_size}
\end{figure}    
\end{frame}


\begin{frame}{Environmental Concerns of E-bike Delivery}
\begin{columns}[T]
    \begin{column}{0.48\textwidth}
        \textbf{E-bike Delivery}
        \begin{enumerate}
            \item \textbf{Battery Life Cycle}
            \begin{itemize}
                \item 2--3 charging cycles/day
                \item Battery lifespan: 1--1.5 years
                \item Lead-acid batteries still common
                \item Poor recycling infrastructure
            \end{itemize}

            \item \textbf{Vehicle Replacement}
            \begin{itemize}
                \item Replaced every 1--2 years
                \item High manufacturing emissions
            \end{itemize}

            \item \textbf{Grid Carbon Intensity}
            \begin{itemize}
                \item China grid $\sim$60\% coal
            \end{itemize}
        \end{enumerate}
    \end{column}

    \begin{column}{0.48\textwidth}
        \begin{enumerate}
            \setcounter{enumi}{3}
            \item \textbf{Operational Inefficiency}
            \begin{itemize}
                \item Road detour: 1.3--1.5$\times$ distance
                \item Stop-and-go traffic
            \end{itemize}

            \item \textbf{Scale}
            \begin{itemize}
                \item 7--10 million riders in China
                \item 40--50 million orders/day
            \end{itemize}
        \end{enumerate}
    \end{column}
\end{columns}
\end{frame}


\begin{frame}{Drone Delivery as Alternative}
\begin{columns}[T]
    \begin{column}{0.48\textwidth}
        \textbf{E-bike Delivery}
        \begin{enumerate}
            \item Battery lifespan: 1--1.5 years
            \item Replaced every 1--2 years
            \item China grid $\sim$60\% coal
            \item Road detour: 1.3--1.5$\times$
            \item 40--50M orders/day
        \end{enumerate}
    \end{column}

    \begin{column}{0.48\textwidth}
        \textbf{Drone Delivery} \textcolor{green!60!black}{(Potential)}
        \begin{enumerate}
            \item Centralized maintenance
            \item Modular component replacement
            \item Same grid, less energy/trip
            \item \textcolor{green!60!black}{Direct line flight (1.0$\times$)}
            \item \textcolor{blue}{$\rightarrow$ Comparative study needed}
        \end{enumerate}
    \end{column}
\end{columns}

\vspace{1em}
\centering
\textbf{Research Question:} Under what conditions does drone delivery produce lower emissions than e-bike delivery?
\end{frame}


\begin{frame}{Review}
    
\end{frame}


\begin{frame}{Contributions}
\begin{itemize}
    \item We develop a \textbf{gravity-based food delivery demand generation method} that produces \textbf{door-to-door order data}, enabling more precise analysis than conventional zone-based or aggregate approaches.

    \item We present the \textbf{first comprehensive comparison} of carbon emissions between \textbf{e-bike and drone delivery systems} across multiple cities in China.

    \item We identify \textbf{key urban features} (e.g., population density, terrain, delivery distance distribution) that determine when drone delivery achieves lower emissions, providing actionable insights for sustainable logistics planning.
\end{itemize}
\end{frame}


\section{Method}


\begin{frame}{OD Generation - Monthly Sales}



\end{frame}


\begin{frame}{Emission Model - E-bike}

The e-bike power model calculates the instantaneous mechanical power required to propel the vehicle at a given velocity \parencite{kirschstein_comparison_2020}.
\[
P_{\text{ebike}} =
P_{\text{rolling}} + P_{\text{air}} + P_{\text{grade}} + P_{\text{acc}}
\]

\begin{description}
  \item[Rolling Resistance Power:]
    \[
    P_{\text{rolling}} = g \cdot C_{\text{roll}} \cdot v \cdot m
    \]

  \item[Aerodynamic Drag Power:]
    \[
    P_{\text{air}} = \tfrac{\rho \cdot v^3}{2 \cdot 10^3} \cdot A \cdot C_{\text{air}} 
    \]

  \item[Grade Climbing Power:]
    \[
    P_{\text{grade}} = g \cdot i \cdot v \cdot m
    \]

  \item[Acceleration Power:]
    \[
    P_{\text{acc}} =
    \frac{0.504 \cdot n_{\text{acc}}}{2 \cdot 10^3 \cdot 3.6} \cdot m \cdot v_{\mathrm{km/h}}^3
    \]
\end{description}
\end{frame}


\begin{frame}{Emission Model - E-bike}

% \footnotesize
\centering
\captionof{table}{Model parameters and default values for e-bike}
\begin{tabular}{c c c c}
\toprule
\textbf{Symbol} & \textbf{Description} & \textbf{Unit} & \textbf{Value} \\
\midrule
$m$ & Total mass & $\mathrm{t}$ & $0.156$ \\
$v$ & Velocity & $\mathrm{m/s}$ & $8.33$ \\
$v_{\mathrm{km/h}}$ & Velocity & $\mathrm{km/h}$ & $30$ \\
$\rho$ & Air density & $\mathrm{kg/m^3}$ & $1.225$ \\
$g$ & Gravitational acceleration & $\mathrm{m/s^2}$ & $9.81$ \\
$A$ & Frontal area & $\mathrm{m^2}$ & $1.0$ \\
$\theta$ & Road grade (slope) & -- & $0.015$ \\
$n_{\text{acc}}$ & Acceleration frequency & $\mathrm{km^{-1}}$ & $1$ \\
$C_{\text{roll}}$ & Rolling resistance coefficient & -- & $0.008$ \\
$C_{\text{air}}$ & Aerodynamic drag coefficient & -- & $0.65$ \\
\bottomrule
\end{tabular}
\end{frame}


\begin{frame}{Emission Model - Drone}

The drone power model calculates the power required for multi-rotor flight at various flight angles \parencite{kirschstein_comparison_2020}.
\[
P_{\text{drone}} 
= P_{\text{induced}} + P_{\text{profile}} + P_{\text{drag}} + P_{\text{climb}} + P_{\text{aux}}
\]
\begin{description}
  \item[Induced Power:]
  \[ P_{\text{induced}} = k_{\text{lift}} \cdot w \cdot T \]
  \item[Profile Power:]
  \[ P_{\text{profile}} = \rho \cdot R \cdot U^3 \cdot \left(1 + 3\left(\frac{v}{U}\right)^2\right) \cdot \frac{\sigma \cdot C_{\text{d,blade}}}{8} \]
  \item[Fuselage Drag Power:]
  \[ P_{\text{drag}} = \frac{1}{2} \rho \cdot C_{\text{air}} \cdot A \cdot v^3 \]
  \item[Climb Power:]
  \[ P_{\text{climb}} = m \cdot g \cdot v \cdot \sin\gamma \]
  \item[Auxiliary Power:]
  \[ P_{\text{aux}} = \text{constant} \]
\end{description}
\end{frame}


\begin{frame}{Emission Model - Drone}

% \footnotesize

\begin{center}
\captionof{table}{Model parameters and default values for drone}
\begin{tabular}{cccc}
\toprule
Symbol & Description & Unit & Default Value \\
\midrule
$m$ & Total mass & $\mathrm{kg}$ & $14.1$ \\ 
$v$ & Flight velocity & $\mathrm{m/s}$ & $16.67$ \\ 
$T$ & Thrust force & $\mathrm{N}$ & $-$ \\ 
$w$ & Induced velocity & $\mathrm{m/s}$ & $-$ \\ 
$U$ & Blade tip speed & $\mathrm{m/s}$ & $-$ \\ 
$\rho$ & Air density & $\mathrm{kg/m^3}$ & $1.225$ \\ 
$g$ & Gravitational acceleration & $\mathrm{m/s^2}$ & $9.81$ \\
$A$ & Fuselage frontal area & $\mathrm{m^2}$ & $0.15$ \\ 
$\gamma$ & Flight path angle & $\mathrm{rad}$ & $-$ \\ 
$k_{\text{lift}}$ & Lift margin factor & $-$ & $1.15$ \\
$C_{\text{air}}$ & Fuselage drag coefficient & $-$ & $0.65$ \\ 
$C_{\text{d,blade}}$ & Blade drag coefficient & $-$ & $0.075$ \\ 
$\sigma$ & Rotor solidity & $-$ & $-$ \\ 
$R$ & Total rotor disc area & $\mathrm{m^2}$ & $-$ \\
$P_{\text{aux}}$ & Auxiliary power & $\mathrm{kW}$ & $0.01$ \\
\bottomrule
\end{tabular}
\end{center}
\end{frame}


\begin{frame}{Emission Model}
\begin{figure}
    \centering
    \includegraphics[width=0.9\linewidth]{../figures/sensitivity/sensitivity_velocity_power_stacked.pdf}
    \caption{\textbf{Power consumption breakdown by component for drone and e-bike at varying velocities.} (a) Drone power components: induced power, profile drag, parasitic drag, climb power, and auxiliary systems. (b) E-bike power components: rolling resistance, aerodynamic drag, grade resistance, and acceleration.}
    \label{fig:sensitivity_power_breakdown}
\end{figure}    
\end{frame}


\section{Results}

\begin{frame}{Sensitivity Analysis}
\begin{figure}
    \centering
    \includegraphics[width=0.7\linewidth]{../figures/sensitivity/sensitivity_distance_route_factor.pdf}
    \caption{\textbf{Two-dimensional sensitivity analysis of delivery distance and route factor on relative emissions.} The heatmap shows the percentage difference in drone emissions relative to e-bike emissions. The black contour line indicates the breakeven point where both delivery modes produce equal emissions. Negative values (green) indicate scenarios where drone delivery is more environmentally favorable.}
    \label{fig:sensitivity_heatmap}
\end{figure}
\end{frame}



\section{References}


\begin{frame}[allowframebreaks]{References}
\printbibliography
\end{frame}


\section{Appendix}

\begin{frame}{Emission Model - Drone}
\begin{center}
\captionof{table}{Model parameters and default values for drone (continued)}
\begin{tabular}{cccc}
\toprule
Symbol & Description & Unit & Default Value \\
\midrule
$r$ & Rotor radius & m & 0.35 \\ 
$\bar{c}$ & Mean blade chord & m & 0.1 \\ 
$n_r$ & Number of rotors & - & 8 \\ 
$n_b$ & Number of blades per rotor & - & 5 \\ 
$R$ & Total rotor disc area & m² & - \\ 
$v_{\text{level}}$ & Horizontal flight velocity & km/h & 60 \\ 
$v_{\text{vertical}}$ & Vertical flight velocity & km/h & 20 \\ 
$D_{\text{body}}$ & Fuselage drag force & N & - \\ 
$h_{\text{takeoff}}$ & Takeoff/landing height & m & 100 \\ 
$t_{\text{hover}}$ & Hover time per delivery & s & 15 \\
$C_l$ & Blade lift coefficient & - & 0.4 \\ 
\bottomrule
\end{tabular}
\end{center}
\end{frame}


\end{document}